\section{Factor Groups}
\textbf{Conventions.} Permutations compose from right to left. The identity is $e$. A group generated by an element $a$ is denoted $\gen{a}$.

\exercise{14.1}{Order of $\Z_6/\gen3$}
The subgroup is $\gen3=\{0,3\}$, of order $2$. The quotient order is the number of cosets, $6/2=3$. Explicitly the cosets are $\{0,3\}$, $\{1,4\}$, $\{2,5\}$.

\exercise{14.2}{Quotient by a product subgroup}
In $\Z_4$, $\gen2$ has order $2$; in $\Z_{12}$, $\gen2$ has order $6$. Their product therefore has order $2\cdot6=12$. The ambient product has order $4\cdot12=48$, so the quotient order is $48/12=4$. All groups here are abelian, so the subgroup is normal and the quotient is defined.

\exercise{14.3}{Quotient by $\gen{(2,1)}$}
In $\Z_4\times\Z_2$, the coordinates $2$ and $1$ each have order $2$, so $(2,1)$ has order $2$ and generates $\{(0,0),(2,1)\}$. The ambient order is $4\cdot2=8$. Thus the quotient has order $8/2=4$.

\exercise{14.4}{Quotient by the second coordinate factor}
The subgroup $\{0\}\times\Z_5$ has five elements in the fifteen-element group $\Z_3\times\Z_5$, so the quotient has order $3$. Equivalently the projection $(a,b)\mapsto a$ has this subgroup as kernel and image $\Z_3$.

\exercise{14.5}{Quotient of $\Z_2	imes\Z_4$ by a diagonal cycle}
The element $(1,1)$ has order $\lcm(2,4)=4$, so its generated subgroup has four elements. The product has $8$ elements, giving a quotient of order $8/4=2$.

\exercise{14.6}{Quotient of $\Z_{12}	imes\Z_{18}$}
The order of $4\in\Z_{12}$ is $12/\gcd(12,4)=3$, and that of $3\in\Z_{18}$ is $18/\gcd(18,3)=6$. Thus $(4,3)$ has order $\lcm(3,6)=6$. Dividing the ambient order $12\cdot18=216$ by this subgroup order gives quotient order $216/6=36$.

\exercise{14.7}{A quotient of $\Z_2	imes S_3$}
The element $(1,\rho_1)$ has order $\lcm(2,3)=6$. Its subgroup equals $\Z_2\times A_3$: both have six elements, and all its powers lie there. This is normal because $A_3\triangleleft S_3$ and the first coordinate is unrestricted. The ambient order is $2\cdot6=12$, so the quotient has order $12/6=2$. Normality matters here because the ambient group is not abelian.

\exercise{14.8}{A diagonal subgroup of a coprime product}
The element $(1,1)$ in $\Z_{11}\times\Z_{15}$ has order $\lcm(11,15)=165$, equal to the ambient order. It therefore generates the whole group, so the quotient has just one element. Its order is $1$.

\exercise{14.9}{Order of $5+\gen4$ in a cyclic quotient}
The subgroup of $\Z_{12}$ is $H=\{0,4,8\}$. The order of $5+H$ is the least positive $k$ with $5k\in H$. Testing multiples modulo $12$ gives $5,10,3,8$ at $k=1,2,3,4$, respectively; the first three are outside $H$ and $8$ is inside. Thus the order is $4$. Equivalently reduce to $5\equiv1$ modulo $4$.

\exercise{14.10}{Order of $26+\gen{12}$}
The quotient $\Z_{60}/\gen{12}$ is identified with $\Z_{12}$ by reduction modulo $12$, because the kernel consists exactly of the multiples of $12$ modulo $60$. The coset of $26$ corresponds to $2\in\Z_{12}$, whose order is $12/\gcd(12,2)=6$. Thus the requested order is $6$.

\exercise{14.11}{Order of a coset in $\Z_3	imes\Z_6$}
Let $H=\gen{(1,1)}$. Its elements are $(t\bmod3,t\bmod6)$. A multiple $k(2,1)$ lies in $H$ if and only if an integer $t$ satisfies
\[t\equiv2k\pmod3,\qquad t\equiv k\pmod6.\]
The second congruence implies $t\equiv k\pmod3$, so consistency is exactly $3\mid k$. This condition is also sufficient by taking $t=k$. Therefore the least positive such $k$ is $3$, giving order $3$.

\exercise{14.12}{Order of a coset in $\Z_4^2$}
The homomorphism $(x,y)\mapsto x-y\pmod4$ has kernel $\gen{(1,1)}$, so the coset of $(3,1)$ corresponds to $3-1=2\in\Z_4$. This element has order $2$. Directly, $(3,1)$ is not on the diagonal, but $2(3,1)=(2,2)$ is, giving the same least positive return.

\exercise{14.13}{Order modulo $\gen{(0,2)}$}
In $\Z_4\times\Z_8$, the subgroup consists of $(0,0),(0,2),(0,4),(0,6)$. For $k(3,1)$ to lie in it, the first coordinate must vanish, requiring $4\mid3k$ or $4\mid k$, and the second coordinate must be even, requiring $2\mid k$. The combined least positive solution is $k=4$. Hence the coset has order $4$.

\exercise{14.14}{Order modulo $\gen{(1,2)}$}
In $\Z_4\times\Z_8$, use $f(x,y)=y-2x\pmod8$. This is well defined: changing $x$ by $4$ changes $-2x$ by $-8$. Its kernel is exactly $\{(t,2t):t\in\Z_4\}=\gen{(1,2)}$. Thus the coset of $(3,3)$ has the same order as $f(3,3)=3-6=-3\equiv5\pmod8$. Since $\gcd(5,8)=1$, its order is $8$. This also verifies that no smaller positive multiple enters the subgroup.

\exercise{14.15}{A coset already equal to the identity}
In $\Z_6\times\Z_8$ one has $2(4,4)=(8,8)=(2,0)$. Hence $(2,0)$ already belongs to $H=\gen{(4,4)}$. The coset $(2,0)+H$ is therefore $H$, the identity of the quotient, and its order is $1$.

\exercise{14.16}{Conjugating a reflection subgroup of $S_3$}
The source convention is $i_g(x)=gxg^{-1}$. Here $\rho_1=(1,2,3)$ and $\mu_1=(2,3)$. Conjugation relabels a transposition by applying the conjugating permutation to its entries:
\[\rho_1(2,3)\rho_1^{-1}=(\rho_1(2),\rho_1(3))=(3,1)=\mu_2.\]
The identity remains the identity, so $i_{\rho_1}[H]=\{\rho_0,\mu_2\}$.

\exercise{14.17}{Correcting the first normal-subgroup definition}
The definition must read: a subgroup $H\le G$ is normal if $gH=Hg$ for every $g\in G$. The original condition $hG=Gh$ for $h\in H$ says nothing about normality, because for any $h\in G$ both $hG$ and $Gh$ are just the whole group $G$. The subgroup being translated must be $H$, and the quantifier must range over all of $G$.

\exercise{14.18}{The conjugation definition of normality}
The stated definition is correct, provided $H$ is a subgroup: $g^{-1}hg\in H$ for every $h\in H$ and $g\in G$ characterizes normality. Although this initially states an inclusion $g^{-1}Hg\subseteq H$, replacing $g$ by $g^{-1}$ gives the reverse inclusion after conjugation, so equality follows. Thus it is equivalent to invariance under every inner automorphism.

\exercise{14.19}{Definition of an automorphism}
An automorphism of $G$ is an \emph{isomorphism} from $G$ onto itself, meaning a bijective homomorphism $G\to G$. A homomorphism merely mapping $G$ into itself is an endomorphism and need not be invertible. For example, the constant map to $e$ is an endomorphism of every group but is not an automorphism unless the group is trivial.

\exercise{14.20}{Why normal subgroups matter}
Normality is exactly the condition making the rule $(aH)(bH)=(ab)H$ independent of the chosen representatives, so that the cosets form the factor group $G/H$. Every normal subgroup is the kernel of the canonical homomorphism $g\mapsto gH$, and conversely every kernel is normal. Thus normal subgroups allow us to collapse specified elements to the identity while retaining a valid group operation, and describe the information lost by homomorphisms.

\exercise{14.21}{Writing a proof about elements of a quotient}
\textbf{(a)} The instructor suspects that the writer is confusing elements of $G$ with elements of $G/H$, which are cosets. Naming quotient elements $a,b$ is formally permissible, but then one cannot invoke commutativity of $G$ on those objects without choosing representatives. The difficulty is not the letter names themselves; it is failure to track which group they belong to.
\textbf{(b)} A clear beginning is: ``Let $aH,bH\in G/H$, where $a,b\in G$. We must show $(aH)(bH)=(bH)(aH)$.''
\textbf{(c)} Since $G$ is abelian, $ab=ba$. Therefore
\[(aH)(bH)=(ab)H=(ba)H=(bH)(aH).\]
The quotient multiplication is valid because $H$ is normal. This proves $G/H$ is abelian.

\exercise{14.22}{A quotient of a torsion group}
\textbf{(a)} As in Exercise 21, the likely error is treating a coset as though it were an element of the original group. The sentence ``let $x\in G/H$'' is not itself false, but the torsion hypothesis is about representatives in $G$, so the proof must display one.
\textbf{(b)} Write: ``Let $aH\in G/H$ with $a\in G$. Since $G$ is torsion, $a$ has finite order $r$.''
\textbf{(c)} Then $(aH)^r=a^rH=eH=H$, the quotient identity. Thus $aH$ has finite order dividing $r$. Since the coset was arbitrary, the quotient is torsion. Its order need not equal the representative's order; it can decrease.

\exercise{14.23}{True or false}
\begin{enumerate}[label=(\alph*)]
\item True: the standard coset product defines a group exactly when $N\triangleleft G$.
\item True: $ghg^{-1}=h$ in an abelian group.
\item True: conjugating $x$ by $g$ gives $gxg^{-1}=x$ when the group is abelian.
\item True: its order is $|G|/|H|$.
\item True: if $a^r=e$, then $(aH)^r=H$.
\item False: $\Z$ is torsion free but $\Z/2\Z$ has a nonidentity element of order $2$.
\item True: $(aH)(bH)=abH=baH=(bH)(aH)$.
\item False: $S_3/A_3\cong\Z_2$ is abelian, although $S_3$ is not.
\item True for positive integer $n$: the class of $1$ generates all $n$ residue classes.
\item False in general: for nonzero integer $n$, $n\R=\R$, since $x=n(x/n)$, so the quotient is trivial, not of order $n$ unless $n=1$.
\end{enumerate}

\exercise{14.24}{The alternating subgroup and its quotient}
For $n\ge2$, the sign map $\operatorname{sgn}:S_n\to\{1,-1\}$ is a surjective homomorphism with kernel $A_n$. A kernel is normal, so $A_n\triangleleft S_n$, and the homomorphism theorem gives
\[S_n/A_n\cong\{1,-1\}\cong\Z_2.\]
Equivalently there are two cosets, the even and the odd permutations; odd times odd is even. For $n=1$, $S_1=A_1$ and the quotient is trivial, which is the small-case exception to the order-$2$ conclusion.

\exercise{14.25}{Well-defined coset multiplication forces an inclusion}
Let $h\in H$. Since $hH=H=eH$, well-definedness of multiplication gives
\[(ha)H=(hH)(aH)=(eH)(aH)=aH.\]
The element $ha$ lies in its own left coset $(ha)H$, so $ha\in aH$. Since this holds for every $h\in H$, it proves $Ha\subseteq aH$. This is the desired missing direction and uses equality of cosets rather than an unproved normality assumption.

\exercise{14.26}{Quotienting out all torsion}
The torsion elements form a subgroup $T$ by the earlier proof for abelian groups, and every subgroup of an abelian group is normal. Suppose $gT$ has finite order in $G/T$. Then $(gT)^n=T$ for some positive $n$, so $g^n\in T$. Because that element is torsion, $(g^n)^m=e$ for some positive $m$. Thus $g^{nm}=e$, so $g\in T$ and $gT=T$. The quotient has no nonidentity torsion elements, which means it is torsion free.

\exercise{14.27}{Conjugacy is an equivalence relation on subgroups}
Reflexivity follows from $H=eHe^{-1}$. For symmetry, if $K=gHg^{-1}$ then $H=g^{-1}Kg$. For transitivity, if $K=gHg^{-1}$ and $L=hKh^{-1}$, then
\[L=hgH g^{-1}h^{-1}=(hg)H(hg)^{-1}.\]
Thus conjugacy of subgroups is reflexive, symmetric, and transitive. Each conjugation maps subgroups to subgroups, since it is an automorphism.

\exercise{14.28}{Normal subgroups as conjugacy classes}
A subgroup $H$ is normal exactly when its conjugacy class among subgroups is the singleton $\{H\}$. Indeed, normality means $gHg^{-1}=H$ for every $g$, so there are no other conjugate subgroups. Conversely a singleton class means every conjugate equals $H$, which is the conjugation criterion for normality. This is a class of subgroups, not a conjugacy class of individual group elements.

\exercise{14.29}{Conjugate reflection subgroups in $S_3$}
The subgroup $\{\rho_0,\mu_2\}$ is generated by the transposition $(1,3)$. Conjugating a transposition gives $(g(1),g(3))$, and as $g$ ranges over $S_3$ this realizes every transposition. Hence the conjugate subgroups are exactly
\[\{\rho_0,\mu_1\},\quad\{\rho_0,\mu_2\},\quad\{\rho_0,\mu_3\}.\]
There are no others, since any conjugate subgroup has two elements and its nonidentity element is a transposition.

\exercise{14.30}{Powers lying in a normal subgroup of finite index}
Interpret $m=[G:H]$ as finite, as required for the integer exponent in the statement. The factor group has order $m$, so Lagrange gives $(aH)^m=H$ for each coset. Since $(aH)^m=a^mH$, this is $a^mH=H$, equivalent to $a^m\in H$. The ambient group itself need not be finite; finite index is enough. If the index is infinite, the displayed power $a^m$ is not an ordinary group power and the statement needs the finite-index qualification.

\exercise{14.31}{Intersections of normal subgroups}
Let $N=\bigcap_{i\in I}N_i$ with each $N_i\triangleleft G$. It is a subgroup: all identities belong, and if $a,b\in N$, then $ab^{-1}$ belongs to every $N_i$, hence to $N$. For $g\in G$ and $x\in N$, normality of every $N_i$ gives $gxg^{-1}\in N_i$ for all $i$, so $gxg^{-1}\in N$. Thus $N$ is normal. For an empty family the intersection is understood to be $G$, which is also normal.

\exercise{14.32}{The normal closure of a subset}
Let $\mathcal N$ be the collection of normal subgroups containing $S$. It is nonempty because $G$ itself belongs to it. Set
\[N_S=\bigcap_{N\in\mathcal N}N.\]
By Exercise 31 this is normal. It contains $S$, since every intersected subgroup does. If $M$ is any normal subgroup containing $S$, then $M\in\mathcal N$ and $N_S\subseteq M$ by the definition of intersection. Thus it is the unique smallest normal subgroup containing $S$, called the normal closure.

\exercise{14.33}{The commutator quotient is abelian}
By definition every commutator $aba^{-1}b^{-1}$ lies in the normal subgroup $C$. Therefore in $G/C$,
\[(aC)(bC)(aC)^{-1}(bC)^{-1}=(aba^{-1}b^{-1})C=C.\]
Multiplying by $(bC)(aC)$ on the right gives $(aC)(bC)=(bC)(aC)$. All quotient elements are such cosets, so the quotient is abelian. This proof uses only the inclusion of all commutators, not that each element of $C$ must itself be a single commutator.

\exercise{14.34}{A unique subgroup of a given order is normal}
For any $g\in G$, conjugation is an automorphism, so $gHg^{-1}$ is a subgroup of the same order as $H$. If $H$ is the only subgroup of that order, $gHg^{-1}=H$ for every $g$. This is exactly normality. Finiteness ensures that ``order'' in the premise is the usual integer cardinality, although the uniqueness argument itself works for arbitrary cardinality.

\exercise{14.35}{Intersecting with a normal subgroup}
For $h\in H$ and $x\in H\cap N$, the conjugate $hxh^{-1}$ lies in $H$ because $H$ is a subgroup, and in $N$ because $N\triangleleft G$. Thus it lies in $H\cap N$, proving $H\cap N\triangleleft H$.
It need not be normal in $G$: take $G=N=S_3$ and $H=\{e,(1,2)\}$. Then $N$ is normal in $G$ and $H\cap N=H$. But conjugating $(1,2)$ by $(1,2,3)$ gives $(2,3)\notin H$, so this intersection is not normal in $G$.

\exercise{14.36}{Intersection of all subgroups of a fixed finite order}
Let $K$ be the intersection of all subgroups of order $s$; the family is nonempty by hypothesis. Conjugation by a fixed $g\in G$ permutes that family, since $H\mapsto gHg^{-1}$ preserves subgroup order and has inverse conjugation by $g^{-1}$. Thus
\[gKg^{-1}=g\left(\bigcap_{|H|=s}H\right)g^{-1}
=\bigcap_{|H|=s}gHg^{-1}=\bigcap_{|H|=s}H=K.\]
Conjugation commutes with intersection because it is a bijection. Since an intersection of subgroups is a subgroup, the equality proves $K\triangleleft G$ even though the individual subgroups need not be normal.

\exercise{14.37}{Automorphisms and inner automorphisms}
\textbf{(a)} The identity is an automorphism. A composite of two bijective homomorphisms is again a bijective homomorphism. The inverse of a bijective homomorphism is a homomorphism: writing $x=\phi(a)$, $y=\phi(b)$ gives $\phi^{-1}(xy)=\phi^{-1}(\phi(ab))=ab=\phi^{-1}(x)\phi^{-1}(y)$. Composition is associative because function composition is. Thus $\Aut(G)$ is a group.
\textbf{(b)} First prove the subgroup requirement: $i_gi_h=i_{gh}$, $i_e$ is the identity, and $(i_g)^{-1}=i_{g^{-1}}$, so $\Inn(G)=\{i_g:g\in G\}$ is a subgroup. For any $\phi\in\Aut(G)$ and $x\in G$,
\[(\phi i_g\phi^{-1})(x)=\phi(g\phi^{-1}(x)g^{-1})=\phi(g)x\phi(g)^{-1}=i_{\phi(g)}(x).\]
Thus conjugation by every automorphism sends inner automorphisms to inner automorphisms, proving $\Inn(G)\triangleleft\Aut(G)$.

\exercise{14.38}{Elements inducing the identity inner automorphism}
The condition $i_g=\mathrm{id}$ means $gxg^{-1}=x$ for every $x\in G$, equivalent to $gx=xg$ for every $x$. Thus the set is the center $Z(G)$. It is normal because it is the kernel of the homomorphism
\[G\to\Aut(G),\qquad g\mapsto i_g,\]
whose homomorphism law is $i_gi_h=i_{gh}$. Alternatively, conjugating a central element leaves it unchanged, directly giving normality.

\exercise{14.39}{An induced map on factor groups}
Define $\phi_*:G/H\to G'/H'$ by $\phi_*(gH)=\phi(g)H'$.
\step{1. Check representatives.} If $gH=kH$, then $k^{-1}g\in H$. By the assumption $\phi[H]\subseteq H'$, its image $\phi(k)^{-1}\phi(g)$ is in $H'$. Thus $\phi(g)H'=\phi(k)H'$ and the proposed map is well defined.
\step{2. Check products.} Normality makes both quotient products valid, and
\[\phi_*((gH)(kH))=\phi(gk)H'=\phi(g)\phi(k)H'
=(\phi(g)H')(\phi(k)H').\]
Therefore $\phi_*$ is a homomorphism. It is natural because it sends a coset to the coset of the original image without any arbitrary choices.

\exercise{14.40}{Normal subgroups defined by determinants}
\textbf{(a)} The determinant is a homomorphism $\det:GL(n,\R)\to\R^*$ because $\det(AB)=\det A\det B$ and invertibility gives nonzero determinant. The determinant-$1$ matrices form its kernel, hence a normal subgroup. Explicitly products and inverses keep determinant $1$, and $\det(BAB^{-1})=\det A$.
\textbf{(b)} The map $A\mapsto|\det A|$ is a homomorphism to $\R^+$, because absolute value is multiplicative. Its kernel is exactly the matrices with determinant $+1$ or $-1$. Therefore this set is also a normal subgroup. The identity belongs to both sets because $\det I_n=1$.

\exercise{14.41}{Multiplication of subsets and the quotient group}
\textbf{(a)} The set product is associative: $(AB)C$ and $A(BC)$ both consist of all $abc$ with $a\in A,b\in B,c\in C$, since multiplication in $G$ is associative. The identity subset is $\{e\}$, because $\{e\}A=A\{e\}=A$. But the powerset is not a group: the empty set has product $\varnothing$ with every subset and cannot have product $\{e\}$ with an inverse.
\textbf{(b)} If $N\triangleleft G$, then for $a,b\in G$,
\[(aN)(bN)=\{an_1bn_2:n_1,n_2\in N\}=abN.\]
For inclusion in $abN$, rewrite $an_1bn_2=ab(b^{-1}n_1b)n_2$, whose final factor lies in $N$. Conversely every $abn$ is the product of $a\in aN$ and $bn\in bN$, so equality holds. Thus cosets are closed under the subset product, and the rule agrees with representative multiplication.
\textbf{(c)} Associativity is inherited from subset multiplication. The subset $N$ is an identity on cosets, since $N(aN)=aN=(aN)N$ by the displayed equality and normality. Each $aN$ has inverse $a^{-1}N$, because their products in both orders equal $N$. Hence the cosets form a group, proved directly from subsets. Its identity is $N$, whereas the identity of the whole powerset monoid is $\{e\}$; they coincide exactly when $N=\{e\}$.
